% Options for packages loaded elsewhere
\PassOptionsToPackage{unicode}{hyperref}
\PassOptionsToPackage{hyphens}{url}
%
\documentclass[
  9pt,
]{article}
\usepackage{amsmath,amssymb}
\usepackage{iftex}
\ifPDFTeX
  \usepackage[T1]{fontenc}
  \usepackage[utf8]{inputenc}
  \usepackage{textcomp} % provide euro and other symbols
\else % if luatex or xetex
  \usepackage{unicode-math} % this also loads fontspec
  \defaultfontfeatures{Scale=MatchLowercase}
  \defaultfontfeatures[\rmfamily]{Ligatures=TeX,Scale=1}
\fi
\usepackage{lmodern}
\ifPDFTeX\else
  % xetex/luatex font selection
  \setmainfont[]{DejaVu Serif}
  \setmonofont[]{DejaVu Sans Mono}
\fi
% Use upquote if available, for straight quotes in verbatim environments
\IfFileExists{upquote.sty}{\usepackage{upquote}}{}
\IfFileExists{microtype.sty}{% use microtype if available
  \usepackage[]{microtype}
  \UseMicrotypeSet[protrusion]{basicmath} % disable protrusion for tt fonts
}{}
\makeatletter
\@ifundefined{KOMAClassName}{% if non-KOMA class
  \IfFileExists{parskip.sty}{%
    \usepackage{parskip}
  }{% else
    \setlength{\parindent}{0pt}
    \setlength{\parskip}{6pt plus 2pt minus 1pt}}
}{% if KOMA class
  \KOMAoptions{parskip=half}}
\makeatother
\usepackage{xcolor}
\usepackage[margin=0.8in]{geometry}
\usepackage{longtable,booktabs,array}
\usepackage{calc} % for calculating minipage widths
% Correct order of tables after \paragraph or \subparagraph
\usepackage{etoolbox}
\makeatletter
\patchcmd\longtable{\par}{\if@noskipsec\mbox{}\fi\par}{}{}
\makeatother
% Allow footnotes in longtable head/foot
\IfFileExists{footnotehyper.sty}{\usepackage{footnotehyper}}{\usepackage{footnote}}
\makesavenoteenv{longtable}
\setlength{\emergencystretch}{3em} % prevent overfull lines
\providecommand{\tightlist}{%
  \setlength{\itemsep}{0pt}\setlength{\parskip}{0pt}}
\setcounter{secnumdepth}{-\maxdimen} % remove section numbering
\ifLuaTeX
  \usepackage{selnolig}  % disable illegal ligatures
\fi
\IfFileExists{bookmark.sty}{\usepackage{bookmark}}{\usepackage{hyperref}}
\IfFileExists{xurl.sty}{\usepackage{xurl}}{} % add URL line breaks if available
\urlstyle{same}
\hypersetup{
  pdftitle={Supplementary Material: Fixed-Point Ridge DCD-RLS Adaptive Filtering for Wearable EEG},
  hidelinks,
  pdfcreator={LaTeX via pandoc}}

\title{Supplementary Material: Fixed-Point Ridge DCD-RLS Adaptive
Filtering for Wearable EEG}
\author{}
\date{}

\begin{document}
\maketitle

All tags refer to files in the code and data archive. stage2 =
\texttt{stage2\_round6\_ridge\_stdout.txt}; stage3 =
\texttt{stage3\_physionet\_stdout.txt}; stage4 =
\texttt{stage4\_stdout.txt}; stage5 = \texttt{stage5\_stdout.txt}.

\hypertarget{s1.-full-methods-with-code-provenance}{%
\section{S1. Full Methods with code
provenance}\label{s1.-full-methods-with-code-provenance}}

Tag key: - \textbf{R4} = \texttt{stage2\_round4\_benchmark.py} -
\textbf{S2c} = \texttt{stage2\_round6\_ridge\_benchmark.py} -
\textbf{S2o} = \texttt{stage2\_round6\_ridge\_stdout.txt} - \textbf{V2}
= \texttt{stage3\_physionet\_benchmark\_auditfix\_v2.py} - \textbf{S3o}
= \texttt{stage3\_physionet\_stdout.txt} - \textbf{S4} =
\texttt{stage4\_physionet\_diagnostics\_auditfix.py} /
\texttt{stage4\_stdout.txt}

\begin{center}\rule{0.5\linewidth}{0.5pt}\end{center}

Throughout, ``pre-specified'' means that hypotheses, metrics and pass
rules were fixed in code before the evaluation runs reported here. Two
deviations are disclosed.

First, the Stage 2 held-out epochs had also been used in earlier
development rounds: self-check T3 reproduces held-out values from those
rounds {[}S2o:16{]}. The Stage 2 test set is therefore not untouched.

Second, in Stage 3 the DCD memory grid was extended from \{5, 7, 9\} to
\{5, 7, 9, 11\}. This followed a pre-declared grid-edge rule, after a
4-record quick run selected the grid top on all four records
{[}V2:479{]}. The 23-record run used the extended grid.

\hypertarget{a.-filter}{%
\subsection{A. Filter}\label{a.-filter}}

\textbf{Stage A} is a first-order DC blocker:

y{[}k{]} = x{[}k{]} − x{[}k−1{]} + (1 − 2\^{}−P\_hp)·y{[}k−1{]}, with
P\_hp = round(log₂(f\_s/(2π f\_c))) {[}V2:60--72{]}.

f\_c = 0.64 Hz gives P\_hp = 7 at 512 Hz {[}S3o:9{]}.

\textbf{Stage B} scales and rounds the primary input d and each
reference channel to Q15. Taps are stacked per channel, L\_ax taps each
{[}V2:260, 295--312{]}. The a-priori output and error are

ŷ = (wᵀu) \textgreater\textgreater{} 15, e = d − ŷ {[}V2:338--340{]}.

R and r are held at Q30. Exponential forgetting λ = 1 − 2\^{}−P is
applied by right shift:

\begin{itemize}
\tightlist
\item
  R ← R − (R \textgreater\textgreater{} P) + uuᵀ;
\item
  r ← r − (r \textgreater\textgreater{} P) + e·u − ((δ
  \textgreater\textgreater{} P)·w) \textgreater\textgreater{} 15, where
  δ = γ·var(u)·2\textsuperscript{P·2}30 {[}V2:326, 350--356{]}.
\end{itemize}

The leakage term implements a persistent ridge penalty δ‖w‖², as in
DCD-RLS with ridge-regression penalties {[}Zakharov \& Nascimento
2013{]}.

The leading-element DCD solver makes at most N\_u = 4 coordinate updates
per sample over M\_b = 16 bit levels, with amplitude bound H = 16
{[}V2:281--283, 360--368{]}. It uses only shifts, additions and
comparisons.

\textbf{Multichannel update.} Because each channel is a tapped delay
line, each sample changes only the M leading columns of R. The remaining
blocks are shifted {[}V2:344--348{]}. This block-shift update was
bit-identical to the full outer-product update {[}S3o:36{]}.

\textbf{Initialisation} differs between stages: - Stage 2 starts from R₀
= max(round(var(u)·2\^{}30), 1)·I {[}S2c:184--185{]}. - Stage 3 starts
from R₀ = 0; the persistent load δ keeps the system solvable
{[}V2:321{]}.

\hypertarget{b.-input-scaling}{%
\subsection{B. Input scaling}\label{b.-input-scaling}}

\begin{itemize}
\tightlist
\item
  \textbf{Stage 2 primary runs} scaled each epoch by its own
  max\textbar x\textbar·1.05 {[}S2c:159, 163{]}. This is not causal.
  Causal validation-set scaling (99.9th percentile × 1.05) changed
  held-out post-switch RRMSE by at most +0.017 {[}S2o:356{]}.
\item
  \textbf{Stage 3} scaled each record by the 99.9th percentile of
  \textbar x\textbar{} over samples {[}1024, first artifact onset), ×
  1.05 × 16 headroom {[}V2:36--37, 295, 307{]}. This skips the first 2
  s. Input clipping was at most 0.061\% of artifact samples
  {[}S3o:39--61{]}.
\end{itemize}

\hypertarget{c.-complexity-and-word-lengths}{%
\subsection{C. Complexity and word
lengths}\label{c.-complexity-and-word-lengths}}

\begin{itemize}
\tightlist
\item
  \textbf{Single channel, L = 16:} plain DCD-RLS uses 48 multiplies and
  at most 228 additions/comparisons per sample; ridge DCD-RLS uses 64
  multiplies and at most 248 {[}S2o:300{]}.
\item
  \textbf{48-tap RefA filter, as implemented:} 288 multiplies per
  sample. That is 48 (output) + 144 (block-shift R) + 48 (e·u) + 48
  (leakage) {[}V2:338, 346, 355--356{]}.
\item
  \textbf{Observed peak signed widths in Stage 2:} R 35 bits; r 31--32
  bits; w 18 bits with ridge vs 22 bits without {[}S2o:297, 299{]}.
\end{itemize}

\hypertarget{d.-stage-2-semi-synthetic-benchmark}{%
\subsection{D. Stage 2: semi-synthetic
benchmark}\label{d.-stage-2-semi-synthetic-benchmark}}

\textbf{Data.} EEGdenoiseNet EEG and EMG epochs {[}Zhang 2021{]} were
upsampled 2× to 512 Hz (1024 samples, 2 s) and z-scored per epoch
{[}R4:76--77, 279--282{]}. Validation sets of 50 EEG and 50 EMG epochs
were disjoint from held-out sets of 100 each, drawn with fixed seeds
{[}S2c:97--103; S2o:12{]}.

\textbf{Artifact.} EMG was passed through a third-order Butterworth
band-pass of 5--100 Hz for the first 512 samples. From sample 512 it was
replaced by −0.8× a 12--75 Hz band-pass {[}R4:290--291; S2c:294--296{]}.
The result was scaled to input SNR ∈ \{−7, −1, +2\} dB. The reference
was the raw EMG plus white noise at 1\%, 10\% or 31.6\% of its RMS
{[}S2c:297--299{]}. With P ∈ \{4, 5\} this gives 18 cells, all with L =
16.

\textbf{Methods compared.}

\begin{longtable}[]{@{}ll@{}}
\toprule\noalign{}
ID & Method \\
\midrule\noalign{}
\endhead
\bottomrule\noalign{}
\endlastfoot
M1 & plain integer DCD-RLS \\
M2 & ridge integer DCD-RLS \\
M3 & float64 ridge exponentially weighted LS \\
M4 & float64 Givens QR-RLS \\
M5 & NLMS \\
M6 & standard float64 RLS \\
\end{longtable}

\textbf{Tuning.} γ ∈ \{0, 2\^{}−8, \ldots, 2\^{}1\} and μ ∈ \{0.05, 0.1,
0.2, 0.3, 0.5, 0.7, 1.0\} were chosen on the validation epochs. The
criterion was the mean of pre-switch RRMSE (samples {[}256, 512)) and
post-switch RRMSE ({[}512, 1024)) {[}S2c:387--388, 411--413{]}.

\textbf{Tests} used held-out post-switch RRMSE, M2 minus the comparator:
- H1 (vs M1), H3 (vs M5) and H4 (vs M4) pass if the mean paired
difference is below 0 and the two-sided Wilcoxon p, Holm-adjusted over
the 18 cells, is below 0.05. - H2 (vs M3) passes if the upper 95\%
bootstrap CI is below 0.02 × mean(M3).

The bootstrap used 10,000 resamples with seed 42 {[}S2c:728--739,
762--783{]}.

\hypertarget{e.-stage-3-real-recordings}{%
\subsection{E. Stage 3: real
recordings}\label{e.-stage-3-real-recordings}}

\textbf{Data.} 23 PhysioNet motion-artifact EEG records {[}Sweeney 2012;
PhysioNet dataset{]} were recorded at 2048 Hz. Channel 1 (EEG1, FPz) is
the reference. Motion was applied at Channel 2 (EEG2, FP1h), 20 mm away
{[}Sweeney 2012{]}. Records were decimated 4× to 512 Hz with polyphase
resampling (padtype = ``line'') {[}V2:204--205, 221--222{]}. Trigger
annotations gave four 60-s artifact intervals per record, 240 s in all
{[}S3o:12--35{]}.

\textbf{References,} all after Stage A, 16 taps per axis:

\begin{longtable}[]{@{}
  >{\raggedright\arraybackslash}p{(\columnwidth - 4\tabcolsep) * \real{0.3333}}
  >{\raggedright\arraybackslash}p{(\columnwidth - 4\tabcolsep) * \real{0.3333}}
  >{\raggedright\arraybackslash}p{(\columnwidth - 4\tabcolsep) * \real{0.3333}}@{}}
\toprule\noalign{}
\begin{minipage}[b]{\linewidth}\raggedright
Reference
\end{minipage} & \begin{minipage}[b]{\linewidth}\raggedright
Built from
\end{minipage} & \begin{minipage}[b]{\linewidth}\raggedright
Taps
\end{minipage} \\
\midrule\noalign{}
\endhead
\bottomrule\noalign{}
\endlastfoot
RefA & ACC\_x1, ACC\_y1, ACC\_z1 & 48 \\
RefB & ACC\_\emph{1 − ACC\_}2 & 48 \\
RefC & RefB plus leaky-integrated velocity and position (pole 1 −
2\^{}−6) & 144 \\
\end{longtable}

RefC's velocity and position were mean-subtracted over the whole record,
so RefC is not strictly causal {[}V2:140--145, 228--235{]}.

\textbf{Fixed settings.} γ = 0.125 was frozen from Stage 2 {[}V2:35{]}.

\textbf{NLMS.} Each reference channel was normalised by its standard
deviation over the calibration window. The update is w ← w + μ·e·u/(uᵀu
+ 10⁻⁴) {[}V2:387, 397{]}.

\textbf{Tuning: leave-one-record-out (LORO).} For each record, P (DCD on
RefA, then reused for RefB and RefC) or μ was the grid value that
maximised the mean of ΔSNR\_blocks + ΔSNR\_clean over the other 22
records {[}V2:458--459, 477--479, 499{]}.

\textbf{Metrics.} ΔSNR is SNR out minus SNR in, relative to EEG1
{[}Sweeney 2012{]}. It was computed separately over artifact blocks and
over clean epochs. Correlation improvement and clean-epoch RRMSE were
also computed {[}V2:419--447{]}.

\textbf{Pre-specified tests.} - P1 tested whether DCD\_RefA exceeded
NLMS\_RefA in median ΔSNR over artifact blocks. The direction was
pre-specified; no numeric pass threshold was coded {[}V2:556{]}. - P2
passes if the upper 95\% CI of the median clean-RRMSE increase of
DCD\_RefA over NONE is below +0.02 {[}V2:558{]}.

CIs are bootstrap intervals of the median (10,000 resamples, seed 42)
{[}V2:462--465{]}, and p-values come from two-sided Wilcoxon tests
{[}V2:553{]}. A sensitivity set of N = 22 excluded eeg\_20, whose
ACC\_x1 channel had 22\% zero-difference samples {[}S3o:32;
V2:534--535{]}.

\hypertarget{f.-exploratory-diagnostics-not-pre-specified-p-unadjusted-s4}{%
\subsection{F. Exploratory diagnostics (not pre-specified; p unadjusted)
{[}S4{]}}\label{f.-exploratory-diagnostics-not-pre-specified-p-unadjusted-s4}}

\begin{itemize}
\tightlist
\item
  \textbf{E1:} N\_u ∈ \{4, 8, 16\} × P ∈ \{7, 9, 11\} on DCD\_RefA.
\item
  \textbf{E2:} integer DCD vs float64 RLS with the same ridge, and vs
  near-unregularised float64 RLS (γ = 10⁻⁶), at P ∈ \{5, 7, 9, 11\}.
\item
  \textbf{E3:} LORO over P ∈ \{5, 7, 9, 11, 13\}.
\item
  \textbf{E4:} trade-off curves over all μ and P on RefA and RefC, with
  post-hoc paired Wilcoxon tests {[}stage5 script{]}.
\item
  \textbf{E5:} LORO selection of P on RefC itself.
\end{itemize}

\begin{center}\rule{0.5\linewidth}{0.5pt}\end{center}

\hypertarget{s1b.-claim-to-code-check}{%
\subsection{S1b. Claim-to-code check}\label{s1b.-claim-to-code-check}}

\begin{longtable}[]{@{}
  >{\raggedright\arraybackslash}p{(\columnwidth - 4\tabcolsep) * \real{0.3333}}
  >{\raggedright\arraybackslash}p{(\columnwidth - 4\tabcolsep) * \real{0.3333}}
  >{\raggedright\arraybackslash}p{(\columnwidth - 4\tabcolsep) * \real{0.3333}}@{}}
\toprule\noalign{}
\begin{minipage}[b]{\linewidth}\raggedright
Claim
\end{minipage} & \begin{minipage}[b]{\linewidth}\raggedright
File:line
\end{minipage} & \begin{minipage}[b]{\linewidth}\raggedright
Code on that line (abridged)
\end{minipage} \\
\midrule\noalign{}
\endhead
\bottomrule\noalign{}
\endlastfoot
×16 headroom, skip 2 s & V2:36--37 & \texttt{HEADROOM\ =\ 16.0},
\texttt{CAL\_START\ =\ 1024} \\
Stage 3 scale & V2:295 &
\texttt{np.percentile(np.abs(d\_sig{[}CAL\_START:cal\_len{]}),\ 99.9))\ *\ 1.05\ *\ HEADROOM} \\
R₀ = 0 (Stage 3) & V2:321 &
\texttt{R\_q30\ =\ np.zeros((L\_tot,\ L\_tot))} \\
R₀ = var·I (Stage 2) & S2c:184--185 &
\texttt{reg\_q30\ =\ max(int(round(delta\_reg\ *\ var\_u\_q15\ *\ Q30)),\ 1)} \\
δ & V2:326 &
\texttt{gamma\ *\ var\_u\ *\ (1\ \textless{}\textless{}\ P\_shift)\ *\ Q30} \\
block shift & V2:346 &
\texttt{R\_lead\_old\ -\ (R\_lead\_old\ \textgreater{}\textgreater{}\ P\_shift)\ +\ np.outer(u\_q15,\ u\_q15{[}lead\_idx{]})} \\
leakage / r update & V2:355--356 &
\texttt{leak\_vec\ =\ ((delta\_q30\ \textgreater{}\textgreater{}\ P\_shift)\ *\ w\_q15)\ \textgreater{}\textgreater{}\ 15} \\
padtype & V2:204 &
\texttt{resample\_poly(...,\ 1,\ 4,\ padtype="line")} \\
RefC non-causal mean & V2:143, 145 & \texttt{vel\ -=\ np.mean(vel)},
\texttt{pos\ -=\ np.mean(pos)} \\
grid extension disclosure & V2:479 &
\texttt{P\_GRID\_DCD\ =\ (5,\ 7,\ 9,\ 11)\ \ \#\ ...\ after\ the\ quick\ run\ picked\ P=9} \\
LORO = mean & V2:459 &
\texttt{np.nanmean({[}scores{[}g{]}{[}j{]}\ ...\ if\ j\ !=\ i{]})} \\
P2 rule & V2:558 &
\texttt{verdict\ =\ "PASS"\ if\ hi\ \textless{}\ 0.02} \\
Stage 2 per-epoch scaling & S2c:159 &
\texttt{np.max(np.abs(d\_sig))\ *\ 1.05} \\
path filters & R4:290--291 &
\texttt{butter(3,\ {[}5/(fs/2),\ 100/(fs/2){]})},
\texttt{butter(3,\ {[}12/(fs/2),\ 75/(fs/2){]})} \\
−0.8 and switch at 512 & S2c:295--296 &
\texttt{n2\ =\ -0.8\ *\ lfilter(...)};
\texttt{np.concatenate({[}n1{[}:512{]},\ n2{[}512:{]}{]})} \\
2× upsampling, z-score & R4:279--282 &
\texttt{resample\_poly(eeg\_raw,\ 2,\ 1,\ axis=1)};
\texttt{R\ /=\ R.std(1)} \\
H1/H3/H4 rule & S2c:785 &
\texttt{mean\_d\ \textless{}\ 0.0\ and\ p\_holm\ \textless{}\ 0.05} \\
H2 rule & S2c:782 & \texttt{ci\_hi\ \textless{}\ margin} (margin =
0.02·mean(M3)) \\
max S3 change & S2o:356 & \texttt{Delta=+0.01708} \\
\end{longtable}

\hypertarget{s2.-stage-2-results-with-line-provenance-stage2-file}{%
\section{S2. Stage 2 results with line provenance (stage2
file)}\label{s2.-stage-2-results-with-line-provenance-stage2-file}}

\begin{longtable}[]{@{}
  >{\raggedright\arraybackslash}p{(\columnwidth - 8\tabcolsep) * \real{0.2000}}
  >{\raggedright\arraybackslash}p{(\columnwidth - 8\tabcolsep) * \real{0.2000}}
  >{\raggedright\arraybackslash}p{(\columnwidth - 8\tabcolsep) * \real{0.2000}}
  >{\raggedright\arraybackslash}p{(\columnwidth - 8\tabcolsep) * \real{0.2000}}
  >{\raggedright\arraybackslash}p{(\columnwidth - 8\tabcolsep) * \real{0.2000}}@{}}
\toprule\noalign{}
\begin{minipage}[b]{\linewidth}\raggedright
H
\end{minipage} & \begin{minipage}[b]{\linewidth}\raggedright
Comparison
\end{minipage} & \begin{minipage}[b]{\linewidth}\raggedright
Cells passed
\end{minipage} & \begin{minipage}[b]{\linewidth}\raggedright
Range of mean difference (M2 − comparator)
\end{minipage} & \begin{minipage}[b]{\linewidth}\raggedright
stage2 lines
\end{minipage} \\
\midrule\noalign{}
\endhead
\bottomrule\noalign{}
\endlastfoot
H1 & M2 \textless{} M1 & 18/18 & −0.146 to −0.048 & 215--234 \\
H2 & M2 non-inferior to M3 (2\% margin) & 8/18 & −0.016 to +0.066 &
235--254 \\
H3 & M2 \textless{} M5 (NLMS μ = 0.2) & 17/18 & −0.601 to −0.020
(failing cell +2 dB, 31.6\%, P = 4; Holm p = 0.233, line 264) &
255--274 \\
H4 & M2 \textless{} M4 (QR-RLS) & 18/18 & −0.338 to −0.092 & 275--294 \\
\end{longtable}

At −7 dB, the ridge raised pre-switch RRMSE relative to plain DCD-RLS in
the 1\% and 10\% noise cells at both P (e.g., 1.198 → 1.322, lines
50--51), but lowered it at 31.6\% noise, P = 4 (2.029 → 1.971, lines
104--105). NLMS had the lowest pre-switch RRMSE in four of the six −7 dB
cells (all P = 4 cells and P = 5 at 31.6\%; lines 56, 83, 110, 191). At
P = 5 with 1\% and 10\% noise, plain DCD-RLS was lowest (1.093, 1.219;
lines 131, 158).

Standard float64 RLS diverged at P = 4 in all 9 cells, with post-switch
RRMSE 401--1268 (lines 58, 67, 76, 85, 94, 103, 112, 121, 130), and was
stable at P = 5 (lines 139--211).

Causal validation-set scaling changed held-out post-switch RRMSE by at
most +0.017 (line 356).

\hypertarget{s3.-full-physionet-results-with-line-provenance}{%
\section{S3. Full PhysioNet results with line
provenance}\label{s3.-full-physionet-results-with-line-provenance}}

On the 23 real motion-artifact EEG recordings from PhysioNet (fs = 512
Hz, L\_ax = 16 taps/axis, Nu = 4, Mb = 16, γ = 0.125 {[}stage3: line
4{]}; 240 s of induced artifact per record across four 60 s intervals
{[}stage3: lines 12--35{]}), the primary pre-specified superiority test
(P1) comparing the integer DCD\_RefA filter (48 total taps, 288
multiplies per sample for the 48-tap RefA filter, as implemented {[}v2
code: run\_multichan\_dcd\_ridge\_int, 48 w·u + 144 block-shift R + 48
e·u + 48 ridge leakage{]}) against NLMS\_RefA on the disturbed
electrode's 3-axis accelerometer reference failed {[}stage3: line
106{]}. Across all 23 records under leave-one-record-out (LORO)
cross-validation (P = 11 selected on 23/23 folds for DCD\_RefA
{[}stage3: line 64{]}; μ = 0.05 selected on 23/23 folds for NLMS\_RefA
{[}stage3: line 62{]}), the paired median difference in artifact-block
SNR change (ΔSNR\_blocks, DCD\_RefA minus NLMS\_RefA) was -2.116 dB
(95\% bootstrap confidence interval {[}-2.535, -1.971{]} dB, two-sided
Wilcoxon signed-rank p = 4.8e-7 {[}stage3: line 106{]}), where DCD\_RefA
had the higher ΔSNR on 1/23 records {[}counted from the per-record
table, stage3: lines 68--90{]}. The pre-specified clean-EEG preservation
test (P2) passed: the paired median clean-epoch RRMSE increase
(DCD\_RefA minus unprocessed NONE) was +0.0145 (95\% CI {[}+0.0107,
+0.0165{]}), with the upper confidence bound below the +0.02
non-inferiority margin {[}stage3: line 107{]}, whereas NLMS\_RefA at μ =
0.05 had a median clean-epoch RRMSE increase of +0.0932 (95\% CI
{[}+0.0667, +0.1053{]}) and failed P2 {[}stage5: line 24{]}. In the N =
22 sensitivity check excluding record eeg\_20, P1 yielded a median
difference of -2.105 dB (95\% CI {[}-2.535, -1.918{]} dB) {[}stage3:
line 127{]} and P2 yielded +0.0146 (95\% CI {[}+0.0107, +0.0170{]})
{[}stage3: line 128{]}, leaving both verdicts unchanged {[}stage3: lines
127--128{]}. All subsequent diagnostic analyses (E1--E5) and post-hoc
comparisons are exploratory (not pre-specified; p unadjusted) {[}stage4:
lines 35--97; stage5: line 4{]}. In E1, raising the coordinate-update
limit Nu from 4 to 16 changed median ΔSNR\_blocks by +0.011 dB at P = 11
{[}stage4: line 353{]}, +0.033 dB at P = 9 {[}stage4: line 268{]}, and
+0.131 dB at P = 7 {[}stage4: line 183{]}, showing that the per-sample
update budget does not limit DCD\_RefA at the selected operating point.
In E2, comparing integer DCD\_RefA against float64 ridge RLS with the
identical ridge load gave median ΔSNR\_blocks differences (integer DCD
minus float64 ridge RLS) of -0.010 dB at P = 11 {[}stage4: line 483{]},
-0.039 dB at P = 9 {[}stage4: line 451{]}, -0.210 dB at P = 7 {[}stage4:
line 419{]}, and -1.150 dB at P = 5 {[}stage4: line 387{]}, while the
ridge load improved upon near-unregularised float64 RLS at every P on
23/23 records across both artifact blocks and clean epochs {[}stage4:
lines 388--389, 420--421, 452--453, 484--485{]}. In E3, extending the
memory grid to P = 13 produced a lower mean LORO score (+0.282 at P = 13
versus +1.043 at P = 11) {[}stage4: lines 491--492{]}, and LORO still
selected P = 11 on 23/23 folds {[}stage4: line 493{]}. Across the E4
trade-off curves on RefA (Fig. 1a; {[}stage4: lines 498--510{]};
{[}stage5: lines 16--26{]}), median ΔSNR\_blocks / ΔSNR\_clean for
DCD\_RefA moved from +4.75 / -5.33 dB at P = 5 {[}stage5: line 16{]}, to
+3.05 / -2.83 dB at P = 7 {[}stage5: line 17{]}, +1.78 / -0.99 dB at P =
9 {[}stage5: line 18{]}, +1.15 / -0.31 dB at P = 11 {[}stage5: line
19{]}, and +0.78 / -0.42 dB at P = 13 {[}stage5: line 20{]}, whereas
NLMS\_RefA across μ in \{0.005, 0.01, 0.02, 0.05, 0.1, 0.2\} spanned
+1.42 / -0.28 dB ({[}stage5: line 21{]}), +1.82 / -0.43 dB ({[}stage5:
line 22{]}), +2.60 / -0.85 dB ({[}stage5: line 23{]}), +3.71 / -1.75 dB
({[}stage5: line 24{]}), +4.39 / -2.81 dB ({[}stage5: line 25{]}), and
+5.10 / -4.20 dB ({[}stage5: line 26{]}). On median values, every
NLMS\_RefA operating point lies on the joint frontier and no DCD\_RefA
point does {[}stage4: lines 498--510{]}; at the two settings that passed
P2 (NLMS μ = 0.005, DCD P = 11), paired differences were not significant
(p = 0.16 artifact blocks, p = 0.111 clean epochs) {[}stage5: lines
5--6{]}. On the 9-axis RefC reference (Fig. 1b; {[}stage4: lines
512--524{]}; {[}stage5: lines 27--37{]}), DCD\_RefC at P = 11 (+3.28 /
-0.76 dB, clean-RRMSE increase +0.0443 {[}+0.0310, +0.0498{]} {[}stage5:
line 30{]}) and P = 13 (+2.35 / -0.74 dB, +0.0367 {[}+0.0184, +0.0526{]}
{[}stage5: line 31{]}) lie on the joint median frontier at the
low-damage end {[}stage4: lines 512--524{]}, though no RefC operating
point passes the +0.02 P2 margin {[}stage5: lines 27--37{]}. Post-hoc
paired Wilcoxon signed-rank tests (two-sided, unadjusted) for DCD\_RefC
at P = 11 minus NLMS\_RefC at μ = 0.005 (+2.65 / -0.84 dB, clean-RRMSE
increase +0.0428 {[}+0.0263, +0.0572{]} {[}stage5: line 32{]}) showed no
statistically significant difference on artifact blocks (+0.0260 dB,
95\% CI {[}-0.1933, +0.2397{]} dB, 12/23 records, p = 0.687 {[}stage5:
line 8{]}), clean epochs (+0.0156 dB, 95\% CI {[}-0.0703, +0.0983{]} dB,
12/23 records, p = 0.622 {[}stage5: line 9{]}), or clean RRMSE (-0.0016,
95\% CI {[}-0.0055, +0.0027{]}, 10/23 records, p = 0.315 {[}stage5: line
10{]}), whereas DCD\_RefC at P = 13 trailed NLMS\_RefC at μ = 0.005 on
artifact blocks by -0.7085 dB (95\% CI {[}-0.8954, -0.3440{]} dB, 4/23
records, p = 0.00915 {[}stage5: line 11{]}) with clean-epoch differences
of +0.0125 dB (95\% CI {[}-0.2581, +0.2204{]} dB, 12/23 records, p =
0.777 {[}stage5: line 12{]}) and -0.0004 RRMSE (95\% CI {[}-0.0113,
+0.0161{]}, 11/23 records, p = 0.964 {[}stage5: line 13{]}). When P was
tuned directly on RefC via LORO (E5), LORO selected P = 5 or 7
{[}stage4: line 1037{]}, and DCD\_RefC minus NLMS\_RefC yielded a median
difference of -1.088 dB over artifact blocks (DCD\_RefC had the higher
ΔSNR on 6/23 records) and -1.073 dB over clean epochs (0/23 records)
{[}stage4: lines 1039--1040{]}.

\hypertarget{s4.-verification-table-every-physionet-number-against-its-source-line}{%
\section{S4. Verification table: every PhysioNet number against its
source
line}\label{s4.-verification-table-every-physionet-number-against-its-source-line}}

\begin{longtable}[]{@{}
  >{\raggedright\arraybackslash}p{(\columnwidth - 4\tabcolsep) * \real{0.4849}}
  >{\raggedright\arraybackslash}p{(\columnwidth - 4\tabcolsep) * \real{0.2199}}
  >{\raggedright\arraybackslash}p{(\columnwidth - 4\tabcolsep) * \real{0.2952}}@{}}
\toprule\noalign{}
\begin{minipage}[b]{\linewidth}\raggedright
\textbf{Number / Value in Text}
\end{minipage} & \begin{minipage}[b]{\linewidth}\raggedright
\textbf{Context in Text}
\end{minipage} & \begin{minipage}[b]{\linewidth}\raggedright
\textbf{Exact Source File \& Line Citation}
\end{minipage} \\
\midrule\noalign{}
\endhead
\bottomrule\noalign{}
\endlastfoot
23 records & Number of PhysioNet EEG recordings evaluated & {[}stage3:
lines 12--35, 106{]}, {[}stage5: line 5{]} \\
fs = 512 Hz, L\_ax = 16, Nu = 4, Mb = 16, γ = 0.125 & Primary DCD\_RefA
configuration and decimated rate & {[}stage3: line 4{]} \\
240 s of artifact in four 60 s intervals per record & Induced artifact
duration per record & {[}stage3: lines 12--35{]} \\
288 multiplies/sample for the 48-tap RefA filter (48 w·u + 144
block-shift R + 48 e·u + 48 ridge leakage) & Per-sample multiply count
of run\_multichan\_dcd\_ridge\_int on 3-axis RefA & {[}v2 code:
run\_multichan\_dcd\_ridge\_int, 48 w·u + 144 block-shift R + 48 e·u +
48 ridge leakage{]} \\
μ = 0.05 on 23/23 folds & Pre-registered LORO selection for NLMS\_RefA &
{[}stage3: line 62{]} \\
P = 11 on 23/23 folds (in P1 sentence) & Pre-registered LORO selection
for DCD\_RefA & {[}stage3: line 64{]} \\
-2.116 dB, {[}-2.535, -1.971{]} dB, p = 4.8e-7 & Pre-registered P1
(DCD\_RefA - NLMS\_RefA ΔSNR\_blocks) & {[}stage3: line 106{]} \\
1/23 records & Records where DCD\_RefA had higher ΔSNR\_blocks than
NLMS\_RefA & {[}counted from the per-record table, stage3: lines
68--90{]} \\
+0.0145, {[}+0.0107, +0.0165{]}, margin +0.02 & Pre-registered P2
(DCD\_RefA - NONE clean RRMSE) & {[}stage3: line 107{]}, {[}stage5: line
19{]} \\
+0.0932, {[}+0.0667, +0.1053{]} & NLMS\_RefA (μ = 0.05) clean-epoch
RRMSE increase vs NONE & {[}stage5: line 24{]} \\
N = 22, -2.105 dB, {[}-2.535, -1.918{]} dB (P1); +0.0146, {[}+0.0107,
+0.0170{]} (P2) & Sensitivity check excluding eeg\_20 & {[}stage3: lines
127--128{]} \\
Lines 35--97 & All stage4 self-checks passed & {[}stage4: lines
35--97{]} \\
Nu = 4 to 16: +0.011 dB (P=11), +0.033 dB (P=9), +0.131 dB (P=7) & E1
update-budget sweep on DCD\_RefA & {[}stage4: lines 183, 268, 353{]} \\
-0.010 dB (P=11), -0.039 dB (P=9), -0.210 dB (P=7), -1.150 dB (P=5) & E2
integer DCD minus float64 ridge RLS (≤ 0.04 dB at P ≥ 9) & {[}stage4:
lines 387, 419, 451, 483{]} \\
23/23 records at P in \{5, 7, 9, 11\} & E2 ridge RLS vs
near-unregularised float64 RLS & {[}stage4: lines 388--389, 420--421,
452--453, 484--485{]} \\
+0.282 (P = 13) vs +1.043 (P = 11); P = 11 on 23/23 folds & E3 grid
extension to P = 13 & {[}stage4: lines 491--493{]} \\
+4.75 / -5.33 dB (P=5), +3.05 / -2.83 dB (P=7), +1.78 / -0.99 dB (P=9),
+1.15 / -0.31 dB (P=11), +0.78 / -0.42 dB (P=13) & E4 DCD\_RefA median
ΔSNR\_blocks / ΔSNR\_clean & {[}stage4: lines 498--510{]}, {[}stage5:
lines 16--20{]} \\
+1.42 / -0.28 dB (μ=0.005), +1.82 / -0.43 dB (μ=0.01), +2.60 / -0.85 dB
(μ=0.02), +3.71 / -1.75 dB (μ=0.05), +4.39 / -2.81 dB (μ=0.1), +5.10 /
-4.20 dB (μ=0.2) & E4 NLMS\_RefA median ΔSNR\_blocks / ΔSNR\_clean &
{[}stage4: lines 498--510{]}, {[}stage5: lines 21--26{]} \\
+0.0121, {[}+0.0099, +0.0161{]} & NLMS\_RefA (μ = 0.005) clean-RRMSE
increase (P2 pass) & {[}stage5: line 21{]} \\
506 rows (lines 527--1033) & E4 CSV block in stage4\_stdout.txt &
{[}stage4: lines 527--1033{]}, {[}stage5: line 1{]} \\
+0.0797 dB, p=0.16 & Post-hoc NLMS\_RefA (μ=0.005) - DCD\_RefA (P=11),
blocks & {[}stage5: line 5{]} \\
+0.0379 dB, p=0.111 & Post-hoc NLMS\_RefA (μ=0.005) - DCD\_RefA (P=11),
clean & {[}stage5: line 6{]} \\
-0.0017, p=0.134 & Post-hoc NLMS\_RefA (μ=0.005) - DCD\_RefA (P=11),
clean RRMSE & {[}stage5: line 7{]} \\
+3.28 / -0.76 dB, +0.0443 {[}+0.0310, +0.0498{]} & DCD\_RefC (P = 11)
medians and clean-RRMSE CI & {[}stage4: lines 512--524{]}, {[}stage5:
line 30{]} \\
+2.35 / -0.74 dB, +0.0367 {[}+0.0184, +0.0526{]} & DCD\_RefC (P = 13)
medians and clean-RRMSE CI & {[}stage4: lines 512--524{]}, {[}stage5:
line 31{]} \\
+2.65 / -0.84 dB, +0.0428 {[}+0.0263, +0.0572{]} & NLMS\_RefC (μ =
0.005) medians and clean-RRMSE CI & {[}stage4: lines 512--524{]},
{[}stage5: line 32{]} \\
+0.0260 dB, p=0.687 & Post-hoc DCD\_RefC (P=11) - NLMS\_RefC (μ=0.005),
blocks & {[}stage5: line 8{]} \\
+0.0156 dB, p=0.622 & Post-hoc DCD\_RefC (P=11) - NLMS\_RefC (μ=0.005),
clean & {[}stage5: line 9{]} \\
-0.0016, p=0.315 & Post-hoc DCD\_RefC (P=11) - NLMS\_RefC (μ=0.005),
clean RRMSE & {[}stage5: line 10{]} \\
-0.7085 dB, p=0.00915 & Post-hoc DCD\_RefC (P=13) - NLMS\_RefC
(μ=0.005), blocks & {[}stage5: line 11{]} \\
+0.0125 dB, p=0.777 & Post-hoc DCD\_RefC (P=13) - NLMS\_RefC (μ=0.005),
clean & {[}stage5: line 12{]} \\
-0.0004, p=0.964 & Post-hoc DCD\_RefC (P=13) - NLMS\_RefC (μ=0.005),
clean RRMSE & {[}stage5: line 13{]} \\
P = 5 or 7; -1.088 dB blocks (6/23); -1.073 dB clean (0/23) & E5 direct
LORO on RefC & {[}stage4: lines 1037, 1039--1040{]} \\
17/18 cells (H3) & Stage 2 semi-synthetic EEGdenoiseNet step-switch
benchmark advantage & {[}stage2\_round6\_ridge\_stdout.txt: H3{]} \\
EEG1 at FPz, EEG2 at FP1h, 20 mm spacing, motion induced at Channel 2 &
Electrode montage and disturbance protocol & Sweeney et al.~(2012, IEEE
TITB 16(5):918--926, doi:10.1109/TITB.2012.2207400) \\
0.414 (median) & Unprocessed (NONE) clean-epoch RRMSE of EEG2 vs EEG1 &
{[}stage4: line 498{]} \\
0.64 Hz (P\_hp = 7 at 512 Hz), -4.17 dB at 0.5 Hz & Stage A high-pass
attenuation & {[}stage3: line 9{]} \\
Below -4.7 dB (-6.48 dB at P=5; -4.72 dB at P=7) & Clean-epoch ΔSNR of
DCD\_RefC at P = 5 and P = 7 & {[}stage4: lines 520--521{]}, {[}stage5:
lines 27--28{]} \\
\end{longtable}

\end{document}
